\chapter{Evaluación y resultados}
\label{cap:resultados}

\section{Metodología}
Se recopilaron 120 proyectos públicos de distinto tipo (artículos de
congreso, artículos de revista, currículos y tesis) y se compiló cada uno
diez veces en una máquina virtual de dos procesadores virtuales y
\SI{4}{GiB} de memoria. Las ejecuciones se dividieron en dos escenarios:

\begin{itemize}
  \item \textbf{Frío}: contenedor nuevo, sin archivos de formato
    pregenerados ni caché de fuentes.
  \item \textbf{Caliente}: imagen con los formatos de los tres motores y la
    base de datos de fuentes de \texttt{luaotfload} generados durante la
    construcción de la imagen.
\end{itemize}

Para cada ejecución se registró el tiempo de cada fase: creación del
contenedor, carga del formato, pasadas del motor, procesamiento de la
bibliografía y escritura de los artefactos.

\section{Latencia}
La tabla~\ref{tab:latencia} resume la latencia de extremo a extremo por tipo
de documento. La creación del contenedor cuesta alrededor de \SI{1,1}{s}
independientemente del documento, y el motor solo domina el tiempo total en
documentos de más de cuarenta páginas.

\begin{table}[htbp]
  \centering
  \caption{Latencia de extremo a extremo en segundos (mediana / percentil 95).}
  \label{tab:latencia}
  \begin{tabular}{@{}lrrr@{}}
    \toprule
    Tipo de documento & Páginas & Frío & Caliente \\
    \midrule
    Artículo de congreso sin bibliografía & 4--6 & 5,8 / 7,9 & 2,9 / 3,6 \\
    Artículo de congreso con \textsc{Bib}\TeX{} & 6--8 & 9,7 / 12,4 & 6,1 / 7,8 \\
    Artículo de revista con biber & 10--14 & 14,2 / 18,0 & 10,5 / 12,9 \\
    Currículo (fontspec, Lua\TeX{}) & 1--2 & 8,3 / 10,1 & 4,0 / 5,2 \\
    Tesis & 80+ & 41,6 / 55,3 & 36,9 / 48,7 \\
    \bottomrule
  \end{tabular}
\end{table}

La figura~\ref{fig:latencia} muestra la latencia mediana en función del
número de páginas para ambos escenarios. La diferencia entre las curvas se
mantiene aproximadamente constante, lo que confirma que las optimizaciones
actúan sobre el coste fijo de arranque y no sobre el tiempo de composición.

\begin{figure}[htbp]
  \centering
  \begin{tikzpicture}
    \begin{axis}[
        width=0.85\textwidth, height=7cm,
        xlabel={páginas}, ylabel={latencia mediana (s)},
        xmin=0, xmax=60, ymin=0, ymax=30,
        grid=major, grid style={gray!25},
        legend pos=north west, legend style={font=\small},
      ]
      \addplot[color=acento, mark=square*, thick]
        coordinates {(2,6.1) (6,7.4) (10,9.0) (20,13.2) (40,21.5) (60,29.0)};
      \addplot[color=principal, mark=*, thick]
        coordinates {(2,3.6) (6,4.7) (10,6.1) (20,10.3) (40,18.4) (60,26.1)};
      \legend{frío, caliente}
    \end{axis}
  \end{tikzpicture}
  \caption{Latencia mediana de compilación según el número de páginas.}
  \label{fig:latencia}
\end{figure}

\section{Comportamiento de la caché}
Durante seis meses de operación, el 68~\% de las peticiones se resolvieron
desde la caché. La tasa de aciertos varió notablemente según el tipo de
cliente: las plataformas docentes que regeneran el PDF de una tarea en cada
visita alcanzaron el 91~\%, mientras que los editores interactivos que
compilan en cada pulsación se mantuvieron por debajo del 20~\%. La inclusión
de la versión de la imagen en la ecuación~\eqref{eq:clave} provocó un vaciado
completo de la caché en cada actualización de \TeX{} Live; la tasa de
aciertos recuperó su valor habitual en unas 36 horas.

\section{Análisis de fallos}
De las compilaciones que no estaban en caché, el 11,4~\% terminó con error.
Los diagnósticos atribuyeron el 54~\% de esos errores a archivos
inexistentes referenciados con \verb|\includegraphics| o \verb|\input|, el
23~\% a secuencias de control no definidas y el 9~\% a paquetes que requieren
la ejecución de programas externos. Solo el 0,3~\% de las compilaciones fue
interrumpido por el tiempo límite y ninguna superó los límites de memoria o
de procesos una vez ajustados los valores de la tabla~\ref{tab:limites}.

\section{Conclusiones}
Los resultados permiten extraer tres conclusiones principales:

\begin{enumerate}
  \item En documentos breves, la sobrecarga de infraestructura pesa más que
    la composición tipográfica; precalentar la imagen elimina la mayor parte
    de ese coste.
  \item Una caché direccionada por contenido resulta muy rentable porque las
    cargas de trabajo reales vuelven a solicitar documentos sin cambios con
    mucha más frecuencia de la que los modifican.
  \item El aislamiento estricto no supuso un obstáculo en la práctica: menos
    del 10~\% de los errores se debieron a funcionalidades prohibidas, y en
    todos los casos existía una alternativa compatible.
\end{enumerate}

Como trabajo futuro se propone la compilación incremental por páginas y la
transmisión de los diagnósticos mientras el motor sigue en ejecución, de
forma que el usuario reciba el primer error antes de que termine la
compilación completa \citep{knuth1986}.
